\section{Track Circuit Branch}
\label{sec:methodology:track-circuit-branch}

The TrackCircuitBranch module enforces automatic safety responses whenever a track segment’s occupancy status changes. When a segment transitions from clear to occupied, the module immediately identifies all associated protecting signals and commands them to display the RED aspect, ensuring that no conflicting train movement can occur. This direct coupling between track detection and signal enforcement embodies one of the most fundamental principles of interlocking: occupied tracks must always be safeguarded by restrictive signals.

The validation process begins by confirming the existence and active status of the occupied segment before retrieving its protecting signals from database metadata and configuration tables. Once identified, these signals are forcefully set to danger and logged for traceability. If the system detects an anomaly—for example, a failure to update one of the signals—the module escalates the response by emitting freeze signals and generating diagnostic entries with timestamps and reasons for failure. This layered enforcement not only guarantees immediate protection but also preserves the forensic trail needed for post-event analysis.

Integration with the broader interlocking engine is seamless. The module is automatically triggered by occupancy monitoring services and backend processes that supervise train movements across the network. Its outputs are reflected simultaneously in the simulation environment, operator interface, and persistent audit layers, ensuring that both real-time operations and long-term compliance requirements are satisfied. For instance, when track segment T22 becomes newly occupied, protecting signals such as S13 and S14 are instantly forced to red. If one fails to respond, the system transitions into a safe freeze state, with detailed logs capturing every step of the enforcement.

Through this design, the TrackCircuitBranch acts as the reactive guardian of the interlocking system, ensuring that the fundamental link between track detection and signal protection is never broken, even in the presence of faults or unexpected conditions.